\begin{claim}{C.21}{Proved under stated hypotheses}{The raw adjoint works for ordinary presheaves over every ring}\label{C21}
Let $A$ be a small ordinary category, let $k$ be any commutative ring, and equip
\[
 \mathcal M=\operatorname{Fun}(A^{\mathrm{op}},\Mod k)
\]
with its pointwise tensor product and constant unit. Assume the manuscript's canonical normalization (c) for $\varphi=\operatorname{colim}_{A^{\mathrm{op}}}$, and assume its full interval condition (d). Then every resolving adjunction $F\dashv G$ obtained from a projective resolution of $\varphi$ in the category of cocontinuous $k$-linear functors induces quasi-inverse equivalences
\[
 \mathcal M[W^{-1}]\simeq D_{\geq0}(k).
\]
In particular the right adjoint in this assertion is the raw $G$, over every coefficient ring, with no finite-component assumption on the chosen resolution. The tensor product and the ordinary-category presentation in this statement are additional restrictions; we do not infer them for an arbitrary small $k$-additive category in the manuscript's corollary.
\end{claim}
\begin{proof}
The zero ring is immediate, so we assume $k\ne0$. We first construct a countable unital subring $s\subseteq k$ such that the inclusion is pure as a map of $s$-modules. Begin with the image $s_0$ of $\mathbb Z$. Given the countable subring $s_j$, consider all finite systems of linear equations with coefficients and constants in $s_j$ that have a solution in $k$. There are only countably many such systems. Choose one solution of each solvable system and let $s_{j+1}$ be the subring generated by $s_j$ and all the finitely many coordinates in these chosen solutions. It is still countable. Set $s=\bigcup_{j\geq0}s_j$.

Every finite linear system with coefficients and constants in $s$ occurs at some stage. If it has a solution in $k$, it consequently has a solution in $s$. This is precisely the finite-equation criterion for the purity of $s\to k$. We recall directly why that criterion gives the required tensor monomorphisms. If an element of an $s$-module $V$ dies in $k\otimes_sV$, express the element using finitely many generators of a free presentation of $V$. The asserted vanishing involves only finitely many relations from that presentation, and is therefore a finite system, with coefficients and constant vector in $s$, expressing that vector as a linear combination of those relations. A solution in $k$ can be replaced by a solution in $s$. The original element was already zero in $V$. Thus
\[
 V\longrightarrow k\otimes_sV
 \quad\text{is monic for every }V\in\Mod s.
 \tag{C-pure-subring}\label{Cpuresubring}
\]
We do not claim that $k$ is flat over $s$; purity of this inclusion is what the proof uses.

We fix explicit bar resolving functors $\mathsf B^r$ for presheaves over each ring $r$. The corresponding projective resolution of the constant covariant $r$-module on $A$ is the functor
\[
 L^r(a)=C_*(N(A/a);r).
\]
We may use unnormalized chains. The overcategory $A/a$ has terminal object $1_a$, so the augmented complex $L^r(a)\to r[0]$ has its usual contraction. In degree $n$ the functor $L^r_n$ is a direct sum of covariant representables, indexed by strings $a_0\to\cdots\to a_n$. Its associated resolving functor is
\[
 \mathsf B^r_n(X)=
 \bigoplus_{a_0\to\cdots\to a_n}X(a_n),
 \qquad
 \mathsf B^r_0(X)=\bigoplus_{a\in A}X(a).
 \tag{C-bar-formula}\label{Cbarformula}
\]
The face maps are those of the bar construction, using the presheaf structure on the last face. These are projective cocontinuous component functors, hence exact. Claim~\ref{C2} identifies $\mathsf B^r$ with $L\operatorname{colim}$.

Let $U$ denote restriction of scalars from $k$ to $s$, on modules, presheaves, or complexes as appropriate. The displayed formula gives natural, differential-compatible identifications
\[
 \mathsf B^s(UY)=U\mathsf B^k(Y),\qquad
 L^k(a)=k\otimes_sL^s(a).
 \tag{C-bar-scalars}\label{Cbarscalars}
\]
The first identity concerns restriction of scalars and requires no flatness.

We now descend the interval to $s$. Regard $J$ as an $s$-module presheaf, and put
\[
 \beta:s_A\longrightarrow k_A\xrightarrow{\alpha}UJ.
\]
For an arbitrary $s$-module presheaf $X$, the map $\beta\otimes_sX$ factors as
\[
 X\longrightarrow k\otimes_sX
 \longrightarrow
 J\otimes_k(k\otimes_sX)
 \cong UJ\otimes_sX.
\]
The first arrow is pointwise monic by \eqref{Cpuresubring}; the second is monic by the original universal interval condition. Moreover, the original condition (d), applied to $k\otimes_sX$, and \eqref{Cbarscalars} give
\[
 \mathsf B^s(UJ\otimes_sX)
 =U\mathsf B^k\bigl(J\otimes_k(k\otimes_sX)\bigr)
 \simeq0.
\]
This proves the whole interval condition over $s$, including degree-zero nullity and all coefficient presheaves.

Canonical normalization descends as well. For any complex $C$ of $s$-modules, purity gives an injection
\[
 H_n(C)\longrightarrow H_n(k\otimes_sC).
 \tag{C-pure-homology}\label{Cpurehomology}
\]
Indeed, if a cycle becomes a boundary after base change, its class in the module $C_n/\operatorname{im}d_{n+1}$ maps to zero after tensoring with $k$. That module injects into its base change by \eqref{Cpuresubring}; the cycle was already a boundary. Apply \eqref{Cpurehomology} to the cone of the canonical augmentation
\[
 \mathsf B^s(s_A)\longrightarrow s[0].
\]
After tensoring with $k$ this is the cone of the normalized augmentation over $k$, hence acyclic by (c). Its homology therefore vanishes already over $s$. Thus (c) holds canonically over $s$.

The other corollary hypotheses are automatic for presheaves on an ordinary small category: the additive envelope of its $s$-linearization has a small set of finitely presented projective generators and projective $s$-module Hom groups; the constant functor is exact and preserves all limits; and the pointwise tensor is cocontinuous. Since $s$ is countable, the pure-global-dimension bound in \citet[Theorem 2.6.3]{BensonGnacadja1999}, applied as in Claim~\ref{C14}, proves that the raw bar right adjoint $\mathsf G^s$ is homotopical and its counit is a quasi-isomorphism on every complex.

The same bar adjunction works raw over $k$. Its right adjoint is explicitly
\[
 \mathsf G^r(K)(a)=
 \operatorname{Hom}_{\Ch_{\geq0}(r)}(L^r(a),K).
\]
Extension--restriction adjunction and \eqref{Cbarscalars} give
\[
 U\mathsf G^k(K)\cong\mathsf G^s(UK),\qquad
 U\mathsf B^k\mathsf G^k(K)
 \cong\mathsf B^s\mathsf G^s(UK).
\]
Under these identifications the counit over $k$ becomes the counit over $s$: both are evaluation of a chain map. Restriction of scalars is exact and conservative. It follows that the raw bar counit over $k$ is a quasi-isomorphism for every $K$, and the bar adjunction gives the asserted equivalence by Claim~\ref{C3}. In particular
\[
 L\varphi\bigl(\mathsf G^k(D^1 V)\bigr)=0
 \quad\text{for every }k\text{-module }V,
 \tag{C-bar-disk-null}\label{Cbardisknull}
\]
without a projectivity requirement on $V$.

It remains to pass from this particular resolution to an arbitrary chosen $F$. We give the retract argument explicitly, since homotopy equivalence of resolutions alone does not automatically settle their raw right adjoints. For $a\in A$, let $R_a$ be the right adjoint of evaluation at $a$. Thus
\[
 R_a(V)(b)=\operatorname{Hom}_k(k[\operatorname{Hom}_A(a,b)],V).
\]
The degree-zero bar formula and the disk adjunction give a natural isomorphism
\[
 \mathsf G^k(D^1V)=\prod_{a\in A}R_a(V).
 \tag{C-bar-dominates}\label{Cbardominates}
\]
Each projective component $F_n$ is a retract of a direct sum of evaluation functors. Its right adjoint $F_n^R$ is consequently a retract of the corresponding product of the $R_a$. For any sequence $(Q_m)_{m\geq1}$ of modules we therefore have
\[
 G\!\left(\prod_{m\geq1}D^mQ_m\right)
 \cong\prod_{m\geq1}F_{m-1}^R(Q_m),
\]
which is a retract of an object of the form $\prod_{t\in T}R_{a_t}(V_t)$ for a set $T$ and modules $V_t$.

Put $V=\prod_{t\in T}V_t$. Since each $R_a$ preserves products, \eqref{Cbardominates} identifies
\[
 \mathsf G^k(D^1V)
 \cong\prod_{(a,t)\in A\times T}R_a(V_t).
\]
The product $\prod_{t\in T}R_{a_t}(V_t)$ is a direct factor of this product: project onto the factors $(a_t,t)$ and use zero maps in every other factor for the section. The assertion also holds when $T$ is empty, since then the object is zero. Thus the chosen $G$ applied to every product of disks is a retract of an object in \eqref{Cbardisknull}, and is itself $L\varphi$-null.

Applying this conclusion to projective disk coefficients proves exactly the necessary-and-sufficient gate of Claim~\ref{C10}. That claim completes the proof for every chosen projective resolving adjunction. The proof has used the ordinary bar resolution and pointwise tensor essentially; it does not identify a general $k$-additive presentation with an ordinary-category presheaf category.
\end{proof}
